\section{Conclusion and Future Work}

This paper introduces \dataname, a tool-agentic dataset containing 1.5M trajectories designed to train better agentic models. We propose a comprehensive pipeline for data generation and demonstrate that models fine-tuned on \dataname achieve superior performance on benchmarks including BFCL-V3 and MCP-Universe.
\dataname represents the first step in a long-term effort to leverage tool use for building stronger LLM agents. Despite being a valuable contribution, we acknowledge our work exhibits certain limitations, which we plan to address through different initiatives.

\textbf{Expanding to More MCP Servers.} While our dataset is comprehensive, it was collected in June 2025, and new servers continue to emerge. We excluded MCP servers requiring special configurations (e.g., requires API keys or account setups), which simplifies the onboarding procedure but may overlook important servers and widely-used scenarios (e.g., Notion and GitHub). Manually onboarding more servers or developing automated onboarding agents could be valuable future work.

\textbf{Expert models to simulate tool-responses.} While real tool execution produces higher-quality results, it is often slow and costly, and therefore, not an option for everyone. To provide an alternative that also yields quality, we plan to develop an expert LLM capable of simulating tool execution. This artificial component will significantly reduce the cost of generating trajectory data involving tool use. Although the idea of tool-execution simulation is known within the community, it has most likely been implemented using off-the-shelf, closed-source LLMs.

\textbf{MCP Benchmark for web search.} As tool-use capabilities become central to both LLMs and LLM-agents, specific scenarios such as web search have gained prominence in the community as a means of synthesizing complex reasoning tasks. To advance this direction, we plan to develop an MCP benchmark focused on web search capabilities.